% AUTHOR'S CLASSROOM HANDOUT: 1.2 (two pages), prepared 2024-01-11.
\section{Completely Symmetric and Antisymmetric States}

Let $\{\ket i\}$ be an orthonormal one-particle basis. An ordered basis vector in the $N$-particle tensor-product space is
\begin{equation}
\ket{i_1,\ldots,i_N}:=\ket{i_1}_1\otimes\cdots\otimes\ket{i_N}_N.
\end{equation}
Write the sum of the permutation operators as
\begin{equation}
\op{\mathcal S}_\pm=\frac1{\sqrt{N!}}
  \sum_{P\in S_N}[\operatorname{sgn}(P)]_\pm\op P,
\end{equation}
where $[\operatorname{sgn}(P)]_+=1$ and $[\operatorname{sgn}(P)]_-=\operatorname{sgn}(P)$. We deliberately use $\mathcal S$ for the unnormalized sum with the square-root prefactor; normalized projectors are introduced below. Changing variables $Q=PP'$ in the group sum yields
\begin{equation}
\op P'\op{\mathcal S}_+=\op{\mathcal S}_+\op P'=\op{\mathcal S}_+,
\qquad
\op P'\op{\mathcal S}_-=\op{\mathcal S}_-\op P'
 =\operatorname{sgn}(P')\op{\mathcal S}_-.
\end{equation}
Consequently $\op{\mathcal S}_+\ket{i_1,\ldots,i_N}$ is symmetric and $\op{\mathcal S}_-\ket{i_1,\ldots,i_N}$ is antisymmetric. They are also orthogonal for $N>1$: an exchange $P_{12}$ preserves the first state and reverses the sign of the second, so their overlap equals its own negative.

\subsection*{1) The symmetrization and antisymmetrization projectors}

Introduce the actual orthogonal projectors
\begin{equation}
\boxed{\op{\mathcal P}_\pm=\frac1{N!}
    \sum_{P\in S_N}[\operatorname{sgn}(P)]_\pm\op P.}
\end{equation}
Their defining properties follow directly from the group multiplication law:
\begin{equation}
\op{\mathcal P}_\pm^\dagger=\op{\mathcal P}_\pm,
\qquad\op{\mathcal P}_\pm^2=\op{\mathcal P}_\pm,
\qquad\op{\mathcal P}_+\op{\mathcal P}_-=0\quad(N>1).
\end{equation}
For example, in the square of the plus projector, for each fixed $R\in S_N$ there are $N!$ pairs $(P,Q)$ with $PQ=R$; therefore
\begin{equation}
\op{\mathcal P}_+^2
=\frac1{(N!)^2}\sum_{P,Q}\op{PQ}
=\frac1{N!}\sum_R\op R=\op{\mathcal P}_+.
\end{equation}
For the minus projector the factor $\operatorname{sgn}(P)\operatorname{sgn}(Q)=\operatorname{sgn}(PQ)$ yields the identical argument. Every completely symmetric state can be written as a linear combination of projected ordered basis states, since the identity resolution can be inserted before projection.

\subsection*{2) Normalization and repeated orbital occupation}

Suppose the ordered list $(i_1,\ldots,i_N)$ contains the orbital $i$ precisely $n_i$ times, with $\sum_i n_i=N$. For bosons, there are $N!/\prod_i n_i!$ distinct permutations of that list, and every distinct ordered product occurs $\prod_i n_i!$ times in the full permutation sum. The norm is therefore
\begin{equation}
\left\|\sum_{P\in S_N}\op P\ket{i_1,\ldots,i_N}\right\|^2
   =N!\prod_i n_i!.
\end{equation}
Thus the normalized Bose state is
\begin{equation}
\boxed{\ket{\{n_i\}}_{\mathrm B}
  =\frac1{\sqrt{N!\prod_i n_i!}}
    \sum_{P\in S_N}\op P\ket{i_1,\ldots,i_N}.}
\end{equation}
In terms of the handout's $\op{\mathcal S}_+$, this is $\ket{\{n_i\}}_{\mathrm B}
=\op{\mathcal S}_+\ket{i_1,\ldots,i_N}/\sqrt{\prod_i n_i!}$. This factorial is essential whenever an orbital is occupied more than once.

For fermions, if $i_a=i_b$ for $a\ne b$, the exchanged ordered product is unchanged but the antisymmetrization sign is reversed, giving zero. Therefore all occupied orbitals must be distinct. For distinct $i_1,\ldots,i_N$, every ordered product in the permutation sum is orthogonal to every other, and
\begin{equation}
\boxed{\ket{i_1,\ldots,i_N}_{\mathrm F}
=\frac1{\sqrt{N!}}\sum_{P\in S_N}
     \operatorname{sgn}(P)\op P\ket{i_1,\ldots,i_N}.}
\end{equation}
In coordinate representation the same expression is the normalized Slater determinant,
\begin{equation}
\Psi_{\mathrm F}(x_1,\ldots,x_N)
=\frac1{\sqrt{N!}}\det\!\begin{pmatrix}
\varphi_{i_1}(x_1)&\cdots&\varphi_{i_N}(x_1)\\
\vdots&\ddots&\vdots\\
\varphi_{i_1}(x_N)&\cdots&\varphi_{i_N}(x_N)
\end{pmatrix}.
\end{equation}
The variable $x$ may include spin as well as position. Both constructions provide complete orthonormal bases of their respective symmetric and antisymmetric $N$-particle sectors. The next chapter assembles sectors with different $N$ into Fock space and derives operator rules directly from these states.